\chapter[From a Parameter Box to Finite Constraints]{From a Parameter Box\\
 to Finite Constraints}
\label{sec:bernstein-motivation}

Parameter-dependent LMIs require a matrix inequality to hold at every point of a continuous parameter domain. The package represents the known data and decision functions on a parameter box with cellwise Bernstein polynomials, then converts the resulting continuum condition into finite constraints for YALMIP. An \term{lpv} induced-\(L_2\)-gain problem provides the first example. The sections that follow develop the general \term{dpd-lmi} and its finite certificate construction. The supporting algebra is developed in Appendices~\ref{app:bernstein-math} and~\ref{app:certificates}, and the API contracts are presented in Chapters~\ref{sec:pdmat}--\ref{sec:pdlmi}.

\section[LPV Induced-L2-Gain Motivation]{Motivation: LPV Induced-\(L_2\)-Gain}

Consider a linear parameter-varying (LPV) plant
\begin{align}
    \dot x(t) & =A(\vect{\rho}(t))x(t)+B_w(\vect{\rho}(t))w(t),\nonumber \\
    z(t)      & =C_z(\vect{\rho}(t))x(t)+D_{zw}(\vect{\rho}(t))w(t),
    \label{eq:lpv-l2-model}
\end{align}
where \(x\in\mathbb R^{n_x}\), \(w\in\mathbb R^{n_w}\), \(z\in\mathbb R^{n_z}\), and the measurable scheduling trajectory satisfies \((\vect{\rho}(t),\dot{\vect{\rho}}(t))\in\mathcal P\times\mathcal R\). For \(\gamma>0\), the zero-state induced-\(L_2\) objective is to certify
\[
    \lVert z\rVert_{\mathcal L_2}<\gamma\lVert w\rVert_{\mathcal L_2}
\]
for every nonzero admissible disturbance and every admissible scheduling trajectory.

A differentiable matrix $P(\vect{\rho})=P(\vect{\rho})^{\mathsf T}$ supplies the scaled bounded-real DPD-LMI \cite{scherer1996}
\begin{align}
    P(\vect{\rho})                                                                                                          & \succ0,
    \label{eq:canonical-positive}                                                                                                         \\
    \begin{bmatrix}
        \dot P(\vect{\rho},\dot{\vect{\rho}})+\operatorname{He}\!\left(P(\vect{\rho})A(\vect{\rho})\right) &
        P(\vect{\rho})B_w(\vect{\rho})                                                                             &
        C_z(\vect{\rho})^{\mathsf T}                                                                                     \\
        B_w(\vect{\rho})^{\mathsf T}P(\vect{\rho})                                                                 &
        -\gamma I_{n_w}                                                                                                    &
        D_{zw}(\vect{\rho})^{\mathsf T}                                                                                  \\
        C_z(\vect{\rho})                                                                                               &
        D_{zw}(\vect{\rho})                                                                                            &
        -\gamma I_{n_z}
    \end{bmatrix} & \prec0
    \label{eq:lpv-spine}
\end{align}
for every \((\vect{\rho},\dot{\vect{\rho}})\in\mathcal P\times\mathcal R\), where \(\operatorname{He}(X)=X+X^{\mathsf T}\) and
\[
    \dot P(\vect{\rho},\dot{\vect{\rho}})
    =\sum_{s=1}^{\ell}\dot\rho_s
    \frac{\partial P}{\partial\rho_s}(\vect{\rho}).
\]
This application is one instance of the general derivative-dependent residual below. The strict signs are analysis notation. The package uses non-strict semidefinite comparisons and applies exactly the numerical margin stated by the user.

\section{General Problem and Assumptions}
\label{sec:general-problem-assumptions}

Let $y=(y_1,\ldots,y_N)$ collect scalar parameter-dependent decision functions. The general task is to satisfy the single residual
\begin{equation}
    \mathcal F(\vect{\rho},\dot{\vect{\rho}};y)
    =F_0(\vect{\rho})
    +\sum_{k=1}^{N}F_k(\vect{\rho})y_k(\vect{\rho})+\sum_{k=1}^{N}\sum_{s=1}^{\ell}
    T_{k,s}(\vect{\rho})\dot\rho_s
    \frac{\partial y_k}{\partial\rho_s}(\vect{\rho})
    \preceq0,
    \quad
    \forall(\vect{\rho},\dot{\vect{\rho}})\in\mathcal P\times\mathcal R.
    \label{eq:general-pdlmi}
\end{equation}
The package fixes the abstract sets and decision class in \cref{eq:general-pdlmi} through the following assumptions.
\begin{itemize}
    \item The scheduling set and rate set are axis-aligned boxes. General polytopic or curved scheduling domains are outside the implemented model:
          \[
              \mathcal P=\prod_{s=1}^{\ell}[\underline\rho_s,\overline\rho_s],
              \qquad
              \mathcal R=\prod_{s=1}^{\ell}[\underline v_s,\overline v_s],
              \qquad
              \mathcal V_{\mathcal R}
              =\prod_{s=1}^{\ell}\{\underline v_s,\overline v_s\}.
          \]
    \item The functions $F_0$, $F_k$, and $T_{k,s}$ are known matrix functions of compatible dimensions, and the $y_k$ are scalar decision functions. A matrix-valued decision is represented by vectorizing its independent scalar entries and assigning one component $y_k$ to each entry, which encodes symmetry through the independent scalar entries. Setting every $T_{k,s}=0$ removes the derivative term and recovers an ordinary derivative-free \term{pd-lmi}.
    \item The theoretical statement uses differentiable $y_k$. The package searches instead over continuous piecewise-Bernstein functions that are differentiable within each physical cell. Their values agree on shared faces, but their one-sided derivatives may differ, so derivative data and the resulting residual conditions remain cell-local at those interfaces.
    \item Known data used by a certificate have cell-local polynomial coefficient evidence, and the assembled residual remains affine in the underlying YALMIP decisions. Supported numeric or known parameter-dependent factors may multiply a decision expression, but products of two decision-bearing expressions are outside the affine model.
    \item Dependence on $\dot{\vect{\rho}}$ is affine. Consequently, checking the distinct vertices in $\mathcal V_{\mathcal R}$ is exactly equivalent to checking the complete rate box $\mathcal R$. There are at most $2^\ell$ such vertices because a direction with equal lower and upper bounds contributes only one value. This exact reduction is separate from the two finite restrictions over the scheduling variables.
    \item MATLAB and YALMIP use non-strict semidefinite comparisons. A theoretical strict condition therefore requires an explicit user-selected margin, for example \(-\mathcal F(\vect{\rho},\dot{\vect{\rho}};y)-\varepsilon I\succeq0\). Failure of a finite certificate leaves feasibility of \cref{eq:general-pdlmi} inconclusive.
\end{itemize}

\section[Gridding the Parameter Domain]{Gridding the Parameter Domain}
\label{sec:intro-gridding}
\index{physical grid}\index{physical node}\index{physical cell}\index{tensor grid}

For each parameter direction \(s=1,\ldots,\ell\), choose an independent, possibly nonuniform axis grid
\[
    \mathcal G_s=\{\rho_s^{(1)},\ldots,\rho_s^{(k_s)}\},
    \qquad
    \underline\rho_s=\rho_s^{(1)}<\cdots<\rho_s^{(k_s)}=\overline\rho_s.
\]
Their Cartesian product is
\[
    \mathcal G=\mathcal G_1\times\cdots\times\mathcal G_\ell,
\]
which contains \(\prod_{s=1}^{\ell}k_s\) physical nodes. The  multi-index \(\vect{c}=(c_1,\ldots,c_\ell)\), with \(c_s\in\{1,\ldots,k_s-1\}\), defines the location of a physical cell $\vect{c}$ and its width along $s$:
\[
    \mathcal H_{\vect{c}}
    =\prod_{s=1}^{\ell}
    [\rho_s^{(c_s)},\rho_s^{(c_s+1)}],
    \qquad
    h_s^{\vect{c}}
    =\rho_s^{(c_s+1)}-\rho_s^{(c_s)}>0.
\]
Thus $\mathcal G$ has $\prod_s k_s$ physical nodes and determines $\prod_s(k_s-1)$ physical cells whose union is exactly $\mathcal P$. The grid exactly partitions the box geometry. Gridding remains a modeling approximation because it restricts the infinite-dimensional search for $y$ to a selected finite-dimensional space of continuous piecewise-polynomial decision functions.

For example, take the following nonuniform two-dimensional grid on $[0,1]^2$:
\[
    \mathcal G_1=\{0,0.4,1\},
    \qquad
    \mathcal G_2=\{0,0.25,1\}.
\]
Their tensor grid \(\mathcal G=\mathcal G_1\times\mathcal G_2\) has nine physical nodes and exactly partitions \([0,1]^2\) into four physical cells. The cell indices are \((1,1)\), \((1,2)\), \((2,1)\), and \((2,2)\).

\begin{figure}[H]
    \centering
    \begin{tikzpicture}[x=8.2cm,y=5.0cm,font=\scriptsize]
        \draw[dp/grid-line,very thick] (0,0) rectangle (1,1);
        \draw[dp/grid-line,very thick] (0.4,0)--(0.4,1);
        \draw[dp/grid-line,very thick] (0,0.25)--(1,0.25);
        \foreach \x/\lab in {0/{0},0.4/{0.4},1/{1}} {
                \node[dp/grid-point] at (\x,0) {};
                \node[dp/grid-point] at (\x,0.25) {};
                \node[dp/grid-point] at (\x,1) {};
                \node[below=5pt] at (\x,0) {$\lab$};
            }
        \foreach \y/\lab in {0/{0},0.25/{0.25},1/{1}} {
                \node[left=5pt] at (0,\y) {$\lab$};
            }
        \node at (0.2,0.125) {\(\vect{c}=(1,1)\)};
        \node at (0.2,0.625) {\(\vect{c}=(1,2)\)};
        \node at (0.7,0.125) {\(\vect{c}=(2,1)\)};
        \node at (0.7,0.625) {\(\vect{c}=(2,2)\)};
        \node[below=14pt] at (0.5,0) {\(\rho_1\)};
        \node[rotate=90,above=14pt] at (0,0.5) {\(\rho_2\)};
    \end{tikzpicture}
    \caption{A nonuniform two-dimensional tensor grid with nine physical nodes and four physical cells. Every point of the box belongs to at least one cell, and the lines are exact cell boundaries of the domain.}
    \label{fig:intro-two-dimensional-grid}
\end{figure}

Each physical cell has its own local coordinate $\vect{\alpha}\in[0,1]^\ell$ and a separate finite coefficient set. Let the direction-wise modeled degree be $\vect{m}=(m_1,\ldots,m_\ell)$ and define the zero-based local-label set $\mathcal I_{\vect{m}}=\prod_{s=1}^{\ell}\{0,\ldots,m_s\}$. The cell then stores $\prod_s(m_s+1)$ coefficients. For the independent choice $\vect{m}=(3,1)$, every two-dimensional cell above has eight labels $\vect{i}=(i_1,i_2)\in\{0,1,2,3\}\times\{0,1\}$. These labels identify in-cell coefficient positions, while the physical sample points remain the grid nodes. They also show that different parameter directions may use different degrees. Shared-face coefficient identification makes the represented decision function continuous across neighboring cells. \Cref{sec:bernstein-tensor-grids} gives the complete traversal convention.

\section[Bernstein Basis on Each Cell]{\texorpdfstring{\hspace{0.3em}Bernstein Basis on Each Cell}{Bernstein Basis on Each Cell}}
\label{sec:continuous-bernstein-representation}
\index{Bernstein basis}\index{cell-local storage}\index{local Bernstein label}

On the unit box $[0,1]^\ell$, the normalized tensor-product Bernstein basis of direction-wise modeled degree $\vect{m}$ is
\[
    B_{\vect{i}}^{\vect{m}}(\vect{\alpha})
    =\prod_{s=1}^{\ell}
    B_{i_s}^{m_s}(\alpha_s)
    =\prod_{s=1}^{\ell}\binom{m_s}{i_s}(1-\alpha_s)^{m_s-i_s}\alpha_s^{i_s},
    \qquad
    \vect{i}\in\mathcal I_{\vect{m}}.
\]

To use this basis on a physical cell, fix any physical cell $\mathcal H_{\vect{c}}$ and map it to the unit box with the inverse of the forward affine map
\[
    \phi_{\vect{c}}:[0,1]^\ell\longrightarrow\mathcal H_{\vect{c}},
    \qquad
    \phi_{\vect{c},s}(\alpha_s)
    =\rho_s^{(c_s)}+h_s^{\vect{c}}\alpha_s.
\]
Thus $\vect{\alpha}=\phi_{\vect{c}}^{-1}(\vect{\rho})$, with $\alpha_s=(\rho_s-\rho_s^{(c_s)})/h_s^{\vect{c}}$. For any cellwise quantity $X$, write its pullback as $X^{(\vect{c})}=X\circ\phi_{\vect{c}}$. A cell-local polynomial is represented by
\begin{equation}
    X^{(\vect{c})}(\vect{\alpha})
    =\sum_{\vect{i}\in\mathcal I_{\vect{m}}}
    B_{\vect{i}}^{\vect{m}}(\vect{\alpha})X^{(\vect{c})}[\vect{i}].
    \label{eq:intro-cell-bernstein}
\end{equation}

The notation is illustrated by the following one-dimensional known-data example, where the degree is still a vector $\vect{m}=(m_1)=(2)$. Consider the scalar polynomial
\[
    a(\rho)=1+\rho+\rho^2
\]
on the two-cell grid $\mathcal G_1=\{0,0.5,1\}$. In this one-dimensional case, $\mathcal H_c=[\rho_1^{(c)},\rho_1^{(c+1)}]$, $h_1^c=\rho_1^{(c+1)}-\rho_1^{(c)}$, and $\phi_c(\alpha)=\rho_1^{(c)}+h_1^c\alpha$, so $\alpha=\phi_c^{-1}(\rho)=(\rho-\rho_1^{(c)})/h_1^c$. Thus $c$ is the one-dimensional physical-cell index, while the Bernstein label remains a separate subscript. The scalar MATLAB input \code{Degree=2} is uniform shorthand for the stored vector $\vect{m}=(2)$ and instructs \code{pdmat} to retain exact function evaluation together with validated Bernstein coefficient data.

\begin{lstlisting}[style=dpmatlab]
grid = {[0 0.5 1]};
a = pdmat(grid, @(rho) 1 + rho + rho.^2, Degree=2);

physicalCells = a.cells()
localLabels = a.lbls()
cell1Coefficients = cell2mat(a.coeffs(1))
cell2Coefficients = cell2mat(a.coeffs(2))
cell1Table = bernTable(a, 1, "oneLine");
disp(cell1Table)
\end{lstlisting}

\begin{lstlisting}[style=dpoutput]
physicalCells = 2×1 double
     1
     2

localLabels = 3×1 double
     0
     1
     2

cell1Coefficients = 1×3 double
    1.0000    1.2500    1.7500

cell2Coefficients = 1×3 double
    1.7500    2.2500    3.0000

    CellSubscript                         Expression
    _____________    _________________________________________________

        {[1]}        "(1-alpha)^2*1 + 2(1-alpha)alpha*1.25 + alpha^2*1.75"
\end{lstlisting}

Fix the first physical cell \(\mathcal H_1=[0,0.5]\). Its local coordinate is \(\alpha^{(1)}(\rho)=2\rho\), and \cref{eq:intro-cell-bernstein} becomes
\[
    a^{(1)}(\alpha)
    =B_0^2(\alpha)\,1
    +B_1^2(\alpha)\,1.25
    +B_2^2(\alpha)\,1.75,
    \qquad \alpha=\alpha^{(1)}(\rho).
\]
This equality represents \(1+\rho+\rho^2\) throughout the complete cell. The three coefficients are basis coordinates inside the cell, and the physical grid points remain the cell endpoints. The upper-face coefficient \(1.75\) of cell 1 equals the lower-face coefficient of cell 2, reflecting continuity of this known polynomial at the shared physical node \(\rho=0.5\). \Cref{sec:bernstein-higher-degree,sec:bernstein-degree-elevation,sec:bernstein-product-math} develop the corresponding higher-degree, elevation, and multiplication rules.

\section{From Cellwise Algebra to Solver Constraints}
\label{sec:cellwise-solver-path}

This section first isolates the Lyapunov block \(\dot P+\operatorname{He}(PA)\) from \cref{eq:lpv-spine} so that the cell, coefficient, and rate-row mechanics remain visible. It then reconnects that block to the complete induced-\(L_2\)-gain model and solver objective. Known matrices are stored as \code{pdmat}, decisions as \code{pdvar}, cellwise rate derivatives are produced by \code{rhodiff}, and finite certificate selections are stored as \code{pdlmi}.

\subsection{Store Known Matrices and Decision Matrices}
\index{pdmat@\texttt{pdmat}}\index{pdvar@\texttt{pdvar}}

The known LPV matrix is affine in \(\rho\) on the two-cell grid \([0,0.5]\cup[0.5,1]\).  Supplying \code{Degree=1} retains exact function evaluation together with cell-local Bernstein coefficient evidence.

\begin{lstlisting}[style=dpmatlab]
yalmip('clear')
grid = {[0 0.5 1]};
A = pdmat(grid, @(rho) [-1, 0.5; -1, -2] + ...
    rho * [-1.3, -20; 2, -10], Degree=1);
P = pdvar(2, grid, "symmetric", Degree=2);
matrixSize = size(P)
physicalCellCount = A.ncell()
knownDegree = A.Degree
decisionDegree = P.Degree
isContinuous = P.IsContinuous
\end{lstlisting}

\begin{lstlisting}[style=dpoutput]
matrixSize =
     2     2
physicalCellCount =
     2
knownDegree =
     1
decisionDegree =
     2
isContinuous =
  logical
   1
\end{lstlisting}

Each cell stores three local symmetric decision matrices.  The upper face of cell 1 and lower face of cell 2 share one YALMIP decision handle, so six local positions contain five global controls.  The interior controls remain independent.

\begin{center}
\begin{tikzpicture}[x=2.05cm,y=1cm,
  face/.style={dp/decision,minimum width=18mm,minimum height=9mm,
    font=\scriptsize,align=center},
  interior/.style={dp/contribution,minimum width=18mm,minimum height=9mm,
    font=\scriptsize,align=center}]
  \node[face] (p0) at (0,0) {$P^{(1)}[0]$\\face};
  \node[interior] (p1) at (1,0) {$P^{(1)}[1]$\\interior};
  \node[face] (p2) at (2,0) {$P^{(1)}[2]$\\shared face};
  \node[interior] (p3) at (3,0) {$P^{(2)}[1]$\\interior};
  \node[face] (p4) at (4,0) {$P^{(2)}[2]$\\face};
  \draw[dp/arrow] (p0)--(p1); \draw[dp/arrow] (p1)--(p2);
  \draw[dp/arrow] (p2)--(p3); \draw[dp/arrow] (p3)--(p4);
  \node[above=5pt of p2,font=\scriptsize,align=center]
    {$P^{(1)}[2]=P^{(2)}[0]$};
  \node[below=7pt of p2,font=\scriptsize,align=center]
    {six local positions, five global matrix controls};
\end{tikzpicture}
\end{center}

\subsection{Multiply in Bernstein Coefficient Space}

On one physical cell \(\mathcal H_{\mathbf c}\), write \(P=\sum_iB_i^mP^{(\mathbf c)}[i]\) and \(A=\sum_jB_j^nA^{(\mathbf c)}[j]\). The ordered product is
\[
PA=\sum_{k=0}^{m+n}B_k^{m+n}R^{(\mathbf c)}[k],\qquad R^{(\mathbf c)}[k]=\sum_{i+j=k}
    \frac{\binom mi\binom nj}{\binom{m+n}{k}}
P^{(\mathbf c)}[i]A^{(\mathbf c)}[j].
\]
For quadratic \(P\) and affine \(A\), the four result coefficients are
\[
R[0]=P_0A_0,\quad
 R[1]=(P_0A_1+2P_1A_0)/3,
\]
\[
R[2]=(2P_1A_1+P_2A_0)/3,\quad R[3]=P_2A_1.
\]
The anti-diagonal grouping \(k=i+j\) and binomial normalization are exact; this is not a pointwise sample product, and matrix order is preserved. \Cref{sec:bernstein-product-math} derives the general tensor-product rule.

\begin{lstlisting}[style=dpmatlab]
PA = P * A;
ATP = A' * P;
degreePA = PA.Degree
degreeATP = ATP.Degree
coefficientCountPA = numel(PA.coeffs(1))
\end{lstlisting}

\begin{lstlisting}[style=dpoutput]
degreePA =
     3
degreeATP =
     3
coefficientCountPA =
     4
\end{lstlisting}

\subsection{Differentiate and Enumerate Rate Vertices}

For a degree-\(m\) cell of width \(h_1^{(\mathbf c)}\), differentiation is the forward difference
\[
  \frac{\partial P}{\partial\rho}
=\frac{m}{h_1^{(\mathbf c)}}\sum_{i=0}^{m-1}
   \bigl(P^{(\mathbf c)}[i+1]-P^{(\mathbf c)}[i]\bigr)B_i^{m-1}(\alpha).
\]
Because the chain-rule term is affine in the bounded rate, checking the two vertices \(-1\) and \(+1\) covers \(\nu\in[-1,1]\). \Cref{sec:bernstein-differentiation} gives the coefficient derivation and its physical-cell scaling.

For \(\ell\) parameters, let \(h_s^{(\mathbf c)}=\rho_s^{(c_s+1)}-\rho_s^{(c_s)}\) be the width of \(\mathcal H_{\mathbf c}\) in direction \(s\), and let \(\mathbf e_s\) be the corresponding unit multi-index. Differentiating the tensor-product representation in \cref{sec:continuous-bernstein-representation} gives
\[
\frac{\partial P^{(\mathbf c)}}{\partial\rho_s}
=\frac{m}{h_s^{(\mathbf c)}}
\sum_{\substack{\mathbf i\in\{0,\ldots,m\}^{\ell}\\ i_s<m}}
\bigl(P^{(\mathbf c)}[\mathbf i+\mathbf e_s]-P^{(\mathbf c)}[\mathbf i]\bigr)
B_{i_s}^{m-1}(\alpha_s)
\prod_{r\ne s}B_{i_r}^{m}(\alpha_r),
\qquad
\dot P^{(\mathbf c)}=\sum_{s=1}^{\ell}\nu_s\frac{\partial P^{(\mathbf c)}}{\partial\rho_s}.
\]
Each partial derivative has degree \(m-1\) in its differentiated direction and degree \(m\) in every other direction. For \(\ell>1\), \code{rhodiff} therefore elevates the shortened direction back to degree \(m\) before summing the rate-weighted partials. If \(\nu_s\in[\underline\nu_s,\overline\nu_s]\), the rate box \(\mathcal R\) has \(2^\ell\) Cartesian vertices \(\boldsymbol\nu^{(q)}\); because \(\dot P\) is affine in \(\boldsymbol\nu\), traversing those vertices is exact. For \(m>0\), each multidimensional cell consequently stores \(2^\ell\) rate rows with \((m+1)^\ell\) tensor-Bernstein coefficients per row.

\begin{lstlisting}[style=dpmatlab]
D = rhodiff(P, [-1 1]);
derivativeDegree = D.Degree
isContinuous = D.IsContinuous
rateRowCount = size(D.coeffs(1),1)
coefficientsPerRow = size(D.coeffs(1),2)
\end{lstlisting}

\begin{lstlisting}[style=dpoutput]
derivativeDegree =
    1
isContinuous =
  logical
   0
rateRowCount =
    2
coefficientsPerRow =
    2
\end{lstlisting}

For example, a degree-one decision on one two-parameter cell has four rate vertices and four tensor-Bernstein labels:

\begin{lstlisting}[style=dpmatlab]
rb2 = [-1 2; -3 5];
P2 = pdvar(1, {[0 2], [10 20]}, Degree=1, RateBounds=rb2);
D2 = rhodiff(P2);
tensorDerivativeDegree = D2.Degree
rateRowCount2D = size(D2.coeffs([1 1]), 1)
coefficientsPerRow2D = size(D2.coeffs([1 1]), 2)
\end{lstlisting}

\begin{lstlisting}[style=dpoutput]
tensorDerivativeDegree =
    1
rateRowCount2D =
    4
coefficientsPerRow2D =
    4
\end{lstlisting}

Thus a coefficient table is indexed independently by physical hypercube, stored rate row, and tensor-Bernstein label. Rate rows are neither physical cells nor additional Bernstein labels. Once the rows are stacked into the finite residual, subsequent certificate formulas suppress the temporary vertex index \(q\).

\subsection{Form a Residual on Every Hypercube}
\index{residual}

The MATLAB residual isolates the Lyapunov block \(\dot P+\operatorname{He}(PA)\) of the full bounded-real LMI in \cref{eq:lpv-spine}; the separate positivity comparison implements the numerical normalization chosen for \cref{eq:canonical-positive}. Addition performs exact degree elevation before the cell-local terms are combined.

\begin{lstlisting}[style=dpmatlab]
residual = D + P * A + A' * P;
Cdecay = residual <= -1e-6 * eye(2);
Cpositive = P >= eye(2);
residualDegreeM = residual.Degree
decayConstraintCount = numel(Cdecay.Constraints)
positivityConstraintCount = numel(Cpositive.Constraints)
\end{lstlisting}

\begin{lstlisting}[style=dpoutput]
residualDegreeM =
    3
decayConstraintCount =
    16
positivityConstraintCount =
    6
\end{lstlisting}

The assembled decay residual has degree \(M=3\): two physical cells times two stored rate rows times four local labels give sixteen coefficient constraints. Positivity has no rate rows and has three degree-two labels per cell. Shared faces are visited from each adjoining cell because certification remains cell-local. For a comparison \(\mathcal F\preceq0\), \code{pdlmi} stores the original residual and sign-normalizes the positive target as \(S=-\mathcal F\); for \(\mathcal F\succeq0\), it uses \(S=\mathcal F\).

\subsection{Choose a Finite Certificate}

Direct Bernstein-coefficient constraints are the default.  Four opt-in selectors rebuild from the same original residual and relation; selections replace rather than compound one another.

\begin{lstlisting}[style=dpmatlab]
Cpolya = Cdecay.applyPolya(1);
Cputinar = Cdecay.applyPutinar();
Csp = Cdecay.applySparseFullBoxPreorder(2);
Cbox = Cdecay.applyFullBoxPreorder();
counts = [numel(Cdecay.Constraints), ...
    numel(Cpolya.Constraints), ...
    numel(Cputinar.Constraints), ...
    numel(Csp.Constraints), ...
    numel(Cbox.Constraints)]
\end{lstlisting}

\begin{lstlisting}[style=dpoutput]
counts =
    16    20    24    24    24
\end{lstlisting}

Direct tests the coefficients of \(S\). P\'olya elevates \(S\) by increment \(d\) before applying the same test. In one parameter, the Gram paths use the parity-specific Markov--Luk\'acs interval representation; for this degree-\(M=3\) residual, absolute order \(r=1\) is the default. Width \(b=2\) therefore reaches the dense endpoint and SparseFullBox canonicalizes to actual FullBox state, while Putinar delegates to the same interval form; all three wrappers store 24 constraints. In multiple parameters, Putinar uses the singleton-generator box module, the sparse full-box tensor-window certificate uses bandwidth \(b\), and the full-box preordering uses every box-generator product. All five software states are finite sufficient certificates for the continuous inequality. Failure of a selected state at fixed \(M,d,r,b\) is inconclusive. \Cref{sec:certificate-direct,sec:certificate-polya,sec:certificate-markov-lukacs,sec:certificate-putinar-schmudgen,sec:certificate-sparse-full-box,sec:certificate-full-box} give the constructions, while \cref{sec:certificate-boundary} states their implemented boundaries.

\subsection{Export to YALMIP and Recover the Solution}
\index{strictness margin}\index{solver diagnostics}

\code{toYalmip} returns exactly the stored finite constraints.  It does not select a solver, add an objective, or insert a strictness margin.

\begin{lstlisting}[style=dpmatlab]
Fdecay = toYalmip(Cdecay);
Fpositive = Cpositive.toYalmip();
opts = sdpsettings('solver', 'sedumi', 'verbose', 0);
sol = optimize([Fdecay, Fpositive], [], opts);
assert(sol.problem == 0, sol.info)
solvedP = value(P);
\end{lstlisting}

The solver choice and \code{sol.problem} semantics belong to YALMIP and the installed solver. After \code{problem == 0}, \code{value(P)} replaces symbolic coefficient payloads with numeric values while preserving the grid, degree, matrix shape, and cell ordering. In the complete bounded-real model, \(\gamma\) is a degree-zero \code{pdvar}; the tested construction and objective extraction appear in \cref{sec:solver-example}. Retain \code{value(gamma.LocalValues\{1\}\{1\})} only after the same successful-status check. Certificate syntax and order boundaries appear in \cref{sec:pdlmi}.

